\section{Preliminaries}\label{sec:prelim}

\subsection{Machines and the statement (T)}\label{sec:prelim-machines}

We use the machine model of the first paper~\cite{FirstPaper2026}, as in the fourth~\cite[Section~2.1]{FourthPaper2026}:
a deterministic Turing machine with a finite state set, a read-only two-way input tape, $k\ge1$ two-way work tapes over a
finite alphabet, and a write-only one-way output tape that receives at most one bit per step. All heads start at
cell~$0$, head moves are $-1$, $0$ or~$+1$, and every step writes the cell under each work head. For $d\ge1$ let
$b(d)=\max(1,\clog d)$. The input is a $d\times d$ matrix with $b(d)$-bit entries written row by row, of length
$L=d^2b(d)$, and the machine must write its transpose in the same encoding. Statement~(T) says that some machine does
this within $c(L+1)$ steps, for a constant~$c$ and all~$d$; its formal version,
\lean{Transpose.LinearTimeTransposable}, is fixed in a frozen file of the first paper's library. Bit~$s$ of entry
$(i,j)$ sits at input position $(di+j)b+s$ and at output position $\tpos((di+j)b+s)=(dj+i)b+s$, so on input~$w$ the
output bit of a transposer at position~$\ell$ is $w_{\tpos(\ell)}$~\cite[Section~2.2 and Lemma~2.5]{FirstPaper2026}.

A machine is \emph{oblivious} at an input length if the positions of all heads at every time, the times at which it
appends output bits, and its halting time do not depend on the input of that length. The \emph{trace} of a run is the
sequence of all head moves together with the append times and the halting time.

\subsection{Walks, hops and owners}\label{sec:prelim-walks}

A \emph{walk} of \emph{length}~$T\ge1$ is a map $x\colon\{0,\dots,T-1\}\to\Z$ with $\abs{x(t+1)-x(t)}\le1$; a step with
$x(t+1)=x(t)$ is a \emph{pause}. The work heads of a machine trace walks $x_1,\dots,x_k$, one per work tape, and these may
be coupled arbitrarily. We follow the fourth paper~\cite[Sections~2.2 and~4.4]{FourthPaper2026}.

For a head~$h$ and a time~$u$, let $\prev_h(u)=\max\{t<u:x_h(t)=x_h(u)\}$ when this set is nonempty. The pair
$(\prev_h(u),u)$ is the \emph{hop} of head~$h$ ending at~$u$, at the \emph{cell} $x_h(u)$; a hop is determined by either
of its endpoints. A hop $(p,u)$ is \emph{long} if $u-p\ge2$. The hops of head~$h$ at a cell~$c$ partition the interval
from the first to the last visit of~$c$ into half-open intervals $(p,u]$, so:
\begin{enumerate}[label=(P\arabic*)]
\item for every time~$\tau$ and every cell~$c$, at most one hop of head~$h$ at~$c$ \emph{spans}~$\tau$, that
is, satisfies $p<\tau\le u$.
\end{enumerate}
Let $(p,u)$ be a long hop at cell~$c$. Then $x_h(t)\ne c$ for $p<t<u$, and since the walk moves by at most one cell per step,
either $x_h(t)>c$ for all $t\in(p,u)$ or $x_h(t)<c$ for all of them. The hop is \emph{upper} in the first case and
\emph{lower} in the second; this is its \emph{side}. The \emph{owner} of the hop is the first time $o\in(p,u)$ at which
$\abs{x_h(o)-c}$ is largest, and $\abs{x_h(o)-c}\ge1$ is its \emph{height}. For a walk without pauses, a \emph{turn} is a
time~$t$ with $x(t+1)-x(t)\ne x(t)-x(t-1)$, and the owner of every long hop is a turn: if the height is~$1$, then
$u=p+2$ and $o=p+1$; otherwise the walk is at distance one less than the height at $o-1$ and at $o+1$.

Two upper long hops of one walk never \emph{cross}: there are no upper long hops $(p,u)$, $(p',u')$ with
$p<p'<u<u'$, and the same holds for lower long hops~\cite[Section~5.1]{FourthPaper2026}. (If $p<p'<u<u'$, then
$x(p')>x(p)$ because $p'\in(p,u)$, while $u\in(p',u')$ and $x(u)=x(p)<x(p')$, contradicting $x>x(p')$ on $(p',u')$.
The argument uses only $x>x(p)$ on $(p,u)$, so it holds with pauses.) Hence two long hops of one walk on the same side
are either disjoint as open intervals or strictly nested, $p<p'<u'<u$; and long hops of one side whose open intervals
contain a common time form a chain under inclusion.

\begin{definition}[Entries]\label{def:entries}
\formal{Owner.lean}{entries}
An \emph{entry} of a cell~$c$ by a walk~$x$ is a time~$t$ with $x(t)=c$ and either $t=0$ or $x(t-1)\ne c$: the first time
of a maximal run of visits.
\end{definition}

\begin{definition}[Owner-load, multiplicity]\label{def:owner-load}
\partformal{TheoremF.lean}{OwnerLoadLE}{the multiplicity.}
Let $H'$ be a set of long hops of walks $x_1,\dots,x_k$ of length~$T$. It has \emph{owner-load at most~$\ell$} (per head
and side) if for every time~$t$, head~$h$ and side~$\sigma$, the hops of~$H'$ of head~$h$ on side~$\sigma$ whose open
intervals contain~$t$ have at most $\ell$ distinct owners. The \emph{multiplicity} $\mu(h,c)$ of a pair $(h,c)$ is the
number of times~$t$ with $x_h(t)=c$; the walk $x_h$ enters~$c$ at most $\mu(h,c)$ times.
\end{definition}

Since $\sum_c\mu(h,c)=T$ for each head, at most $kT/(V+1)$ pairs $(h,c)$ have multiplicity greater than~$V$, and the
same bound holds for pairs with more than $V$ entries.

\subsection{Walk DAGs, cuts and labels}\label{sec:prelim-dag}

The \emph{walk DAG} of $x_1,\dots,x_k$~\cite[Section~4.4]{FourthPaper2026} has vertex set $[0,T)=\{0,\dots,T-1\}$ and two
kinds of edges: \emph{state edges} $t\to t-1$ for $1\le t<T$, and \emph{hop edges} $u\to p$, one for every long hop
$(p,u)$ of any head. A pause joins the same two vertices as a state edge and is identified with it, and an edge is
determined by its endpoints: a pair that is a long hop of several heads is a single edge. For a set $H'$ of long hops,
$G(H')$ denotes the DAG of the state edges and of the edges $u\to p$ with $(p,u)\in H'$; for a set~$R$ of pairs
$(h,c)$, $G-R$ denotes $G(H')$ for the set $H'$ of long hops at pairs outside~$R$. A \emph{cut} is a set of edges; it has
\emph{depth at most~$D$} if every path that uses no edge of the cut contains at most $D$ hop edges.

A \emph{label} with $q$ values is a map $\Psi\colon[0,T)\to\{0,\dots,q-1\}$. A hop edge $u\to p$ is \emph{violated} if
$\Psi(u)\le\Psi(p)$, a time $t\ge1$ is a \emph{descent} if $\Psi(t)<\Psi(t-1)$, and $F_\Psi$ is the cut consisting of
the violated hop edges and of the state edges $t\to t-1$ at descents. Its size is the \emph{cost} of~$\Psi$.

\begin{lemma}[Cuts are labels]\label{lem:labels}
\partformal{Labels.lean}{depth_le_of_label}{the last sentence, the consequence for walk DAGs.}
Let $\Psi$ be a map from a set to $\{0,\dots,q-1\}$, and let $f(0),f(1),\dots,f(n)$ be a sequence of points of that set
with some of its steps $i<n$ marked. If $\Psi(f(i+1))\le\Psi(f(i))$ for every $i<n$, and $\Psi(f(i+1))<\Psi(f(i))$ for
every marked~$i$, then at most $q-1$ steps are marked. Consequently, for every label $\Psi$ with $q$ values of a
walk DAG, or of a sub-DAG $G(H')$, the cut $F_\Psi$ has depth at most $q-1$.
\end{lemma}

\begin{outline}[Outline of the formal proof of the statement about sequences]
By induction on~$n$, the number of marked steps plus $\Psi(f(n))$ is at most $\Psi(f(0))$
(\lean{card_marked_add_le}), and $\Psi(f(0))\le q-1$.
\end{outline}

\begin{proof}[Proof of the last sentence (not formalized)]
Let $t_0\to t_1\to\dots\to t_n$ be a path that avoids $F_\Psi$, and mark its hop edges. Along a state edge
$t_i\to t_{i+1}=t_i-1$, the time $t_i$ is not a descent, so $\Psi(t_{i+1})\le\Psi(t_i)$; along a hop edge, which is not
violated, $\Psi(t_{i+1})<\Psi(t_i)$. By the statement about sequences, the path has at most $q-1$ hop edges.
\end{proof}

The lemma is one direction of the fourth paper's lemma ``cuts are labels''~\cite[Lemma~5.2]{FourthPaper2026}, which also
proves the converse: for every~$D$, the smallest size of a cut of depth at most~$D$ equals the smallest cost of a label
with $D+1$ values. All cuts in this paper are built from labels.

\subsection{Depth reduction and the crux}\label{sec:prelim-tcp}

\begin{definition}[TCP and depth-robustness]\label{def:tcp}
\emph{TCP holds} for a family of $k$-tuples of walks if for every $\eps>0$ there is an integer~$D$ such that the walk DAG
of every member, of length~$T$, has a cut of at most $\eps T$ edges and depth at most~$D$. $\TCP(k)$ is the statement
that TCP holds for the family of all $k$-tuples of walks~\cite[Problem~5.1]{FourthPaper2026}. A family is
\emph{depth-robust} if TCP fails for it: for some $\eps>0$ no single~$D$ serves all members.
\end{definition}

For an oblivious machine~$M$ at length~$L$ with halting time~$T$, the \emph{dependency DAG}
$H_M$~\cite[Definition~4.1]{FourthPaper2026} has vertex set $[0,T)$, state edges $t\to t-1$, and cell edges
$t\to\prev_j(t)$ along the work-head walks. Every edge of $H_M$ joins the endpoints of an edge of the walk DAG of the
work heads: a cell edge of length~$1$ is a pause, and a cell edge of length at least~$2$ is a long hop. The input head
contributes no edges. The same definition applies to a single trace~$\tau$ of any machine at length~$L$, with the head walks,
append times and halting time of~$\tau$; we write $H_\tau$ for this \emph{dependency DAG of the trace}. All inputs that
follow~$\tau$ share this schedule, and the fourth paper's long-edge conversion~\cite[Lemma~4.21]{FourthPaper2026}, whose
proof uses only the DAG, applies to~$H_\tau$. A \emph{genie} is a set $\Gamma\subseteq[0,T)$ of revealed steps; it reveals the values
$v_s$, $s\in\Gamma$ (state after the step, symbols written, output bit), with $\beta\ge\beta(M)$ bits per value, at
\emph{cost} $\beta\card\Gamma$. A path $t_0\to\dots\to t_m$ of $H_M$ is $\Gamma$-free if $t_1,\dots,t_m\notin\Gamma$,
and $\Gamma$ has \emph{full depth at most~$D$} if every $\Gamma$-free path has at most $D$ edges. The cone lemma
\cite[Lemma~4.4 and Corollary~4.5]{FourthPaper2026}, with written proofs there, says that output bit~$\ell$, appended at step~$t_\ell$, is a
function of the revealed word $G_\Gamma(w)=(v_s(w))_{s\in\Gamma}$ and of the input bits at the positions $I_\Gamma(\ell)$
read at the ends of $\Gamma$-free paths from~$t_\ell$, and that $\card{I_\Gamma(\ell)}\le(k+1)^{D+1}$ if $\Gamma$ has
full depth at most~$D$.

The fourth paper calls \emph{depth-TCP} the statement that genies of cost $\eps L$ and bounded full depth exist for all
oblivious machines with $k$ work tapes running in time $C(L+1)$~\cite[Definition~4.7]{FourthPaper2026}. A
\emph{side-information system} at size~$d$ is a family $I=(I(\ell))_{\ell\in[0,L)}$ of subsets of $[0,L)$, and an
\emph{index code} for~$I$ with $m$ bits is a map $G\colon\{0,1\}^L\to\{0,1\}^m$ such that $G(w)$ and $w|_{I(\ell)}$
determine $w_{\tpos(\ell)}$. For a class $\mathcal S$ of systems and $\rho\ge0$, $\Crux(\mathcal S,\rho)$ says that
for every $\eta>0$ there is $d_0$ such that every index code for every system of $\mathcal S$ at a size $d\ge d_0$ has at
least $(\rho-\eta)L$ bits; $\mathcal S_{k,C,\beta,D}$ is the class of the systems $I_\Gamma$ of oblivious machines with $k$
work tapes, value width at most~$\beta$ and halting time at most $C(L+1)$ that append exactly $L$ bits, for genies
of full depth at most~$D$~\cite[Definition~4.8]{FourthPaper2026}. Every fibre of an index code for~$I$ is a binary storage
code for the relation $N_I(p)=I(\tpos(p))$, in which every symbol is a function of the symbols in its recovery set
\cite[Remark~4.9(b)]{FourthPaper2026}.

\begin{theorem}[Theorem D of the fourth paper]\label{thm:D4}
\cited{the fourth paper~\cite[Theorem~4.10]{FourthPaper2026}, with a complete written proof there; not formalized.}
If $\mathrm{depth}\text{-}\TCP(k,C,\beta,\eps,D)$ and $\Crux(\mathcal S_{k,C,\beta,D},\rho)$ hold with $0<\eps<\rho$,
then for every machine with $k$ work tapes and value width at most~$\beta$ there is $d_0$ such that for no $d\ge d_0$ is it
oblivious at length $L=d^2b(d)$ with halting time at most $C(L+1)$ and a transposer of all $d\times d$ matrices.
\end{theorem}

\paragraph{Why it matters.} Theorem~D splits a proof that no oblivious machine transposes in linear time into a
structural half, depth-TCP, and a coding half, the crux. \Cref{sec:nesting,sec:onehead,sec:dnp} work on the structural half through cuts of
walk DAGs, \cref{sec:linear,sec:crux} on the coding half, and \cref{sec:transposers} combines the two. An oblivious
linear-time transposer must violate, for its own parameters, one of the two hypotheses.

The structural hypothesis is a statement about walk DAGs. Every application below converts a cut of the walk DAG into a
genie, as in the proof of the fourth paper's Corollary~B.

\begin{lemma}[From cuts to genies]\label{lem:cut-genie}
\notformalized
Let $M$ have $k$ work tapes and value width at most~$\beta$, let $C\ge1$ and $0<\eps\le1$, and put
$P=\lceil4\beta C/\eps\rceil$. Suppose that at length~$L\ge1$ the run of $M$ along a trace of halting time
$T\le C(L+1)$ is given, and that the walk DAG of its work-head walks on $[0,T)$ has a cut of at most
$\eps T/(4\beta C)$ edges and depth less than~$D_1$. Then the dependency DAG of that trace has a genie of cost at most
$\eps L$ and full depth at most $D_1P-1$.
\end{lemma}

\begin{informalproof}
Let $F$ be the cut and $\Gamma_0$ the set of steps~$s$ at which some edge $t\to s$ of~$F$ ends, so
$\card{\Gamma_0}\le\card F$. Every edge of a $\Gamma_0$-free path of the dependency DAG ends outside~$\Gamma_0$, so its
endpoints are not those of an edge of~$F$; read in the walk DAG, the path avoids~$F$, and its edges of length at least~$2$
are long hops. Hence every $\Gamma_0$-free path has at most $D_1-1$ edges of length greater than~$1$. By the long-edge
conversion of the fourth paper~\cite[Lemma~4.21]{FourthPaper2026} (a written proof there), with $\mu=1<P$ and
$D'=D_1-1$, the genie
$\Gamma=\Gamma_0\cup\{t<T:t\equiv P-1\pmod P\}$ has full depth at most $D_1P-1$ and
$\card\Gamma\le\card F+\lfloor T/P\rfloor$. Its cost is at most
$\eps T/(4C)+\beta T/P\le\eps T/(2C)\le\eps(L+1)/2\le\eps L$.
\end{informalproof}

The lemma uses the trace only through its dependency DAG. When the machine is oblivious and transposes, the revealed word
of the genie is an index code with at most $\eps L$ bits for the system $I_\Gamma\in\mathcal S_{k,C,\beta,D_1P}$, by the
cone lemma and correctness, exactly as in the proof of Theorem~D.

\subsection{The deep-nesting pair}\label{sec:prelim-dnp}

We recall the instance that the fourth paper regards as the sharpest open case of $\TCP(2)$, with the parameter $g=1$
throughout. Let $s_0>m\ge0$ be integers, $\Lambda=2^{s_0}$ and $R_i=2^{s_0-i}$ for $0\le i\le m$, so $R_0=\Lambda$
and $R_m\ge2$. A \emph{level-$m$ stroke} of direction $\sigma\in\{-1,1\}$ is a run of $R_m$ steps equal to~$\sigma$. For
$i<m$, a \emph{level-$i$ stroke} of direction~$\sigma$ is a \emph{segment} of $R_i$ steps equal to~$\sigma$ followed by a
\emph{V} of level $i+1$: a level-$(i+1)$ stroke of direction~$-\sigma$ (the \emph{back stroke}) followed by a
level-$(i+1)$ stroke of direction~$\sigma$ (the \emph{forward stroke}). The walk $U=U(s_0,m,1)$ satisfies $U(0)=0$,
is periodic with period
\[
  P=4(m+1)\Lambda,
\]
and its steps over one period are those of four level-$0$ strokes of directions $1,-1,-1,1$
\cite[Definition~5.3]{FourthPaper2026}. A level-$i$ stroke has $n_iR_i$ steps, where $n_i=m-i+1$, and moves the head by
$\sigma R_i$; U has no pauses. The first time of a back stroke is the \emph{apex} time of its V, the first time of its
forward stroke the \emph{valley} time; the V is \emph{attached} at its \emph{base}, the cell where it starts and ends, and
its \emph{valley cell} is at distance $R_i$ from the base, for a V of level~$i$. The turns of~$U$ are exactly the apex and
valley times of all V's, each of the level of its V, and the first times of the second and fourth level-$0$ strokes of
every period, of level~$0$, at the \emph{tips} $\pm\Lambda$~\cite[Lemma~5.4]{FourthPaper2026}, with a written proof
there. A period contains $2^{i+1}$ V's of level~$i$ for $1\le i\le m$, as counted in the proof of the fourth paper's
Proposition~5.13, hence $2^{i+2}$ turns of level~$i$ and two tips; all turn times and turn positions are multiples
of~$R_m$~\cite[Remark~5.23]{FourthPaper2026}.

For an integer shift~$\xi$ the pair is $x_1=U$ and $x_2(t)=x_1(t+\xi)$; the \emph{deep-nesting pair}
$\DNP(s_0,m,1)$ of the fourth paper is the pair for a particular odd shift, the golden shift
\cite[Definition~5.9]{FourthPaper2026}, and we call the pair for any odd shift a deep-nesting pair. A pair $(p,u)$ is a hop
of head~$2$ if and only if $(p+\xi,u+\xi)$ is a hop of head~$1$. For $r\ge1$ the \emph{$r$-fold cover} is the pair on the
cyclic time axis $\Z_T$, $T=rP$; its hops are the residues of the hops on~$\Z$, every time is the earlier and the
later endpoint of exactly one hop of each head, and there are $n_{\mathrm h}=2T$ hops in all. Each head visits every cell of
$[-\Lambda,\Lambda]$ in every period, so every time lies in the half-open interval $(p,u]$ of exactly $2(2\Lambda+1)$ hops of
the two heads, one per head and cell. A \emph{label} on the cover is a map $G\colon\Z_T\to\{0,\dots,q-1\}$; a hop $(p,u)$
is violated if $G(u)\le G(p)$, a time~$t$ is a descent if $G(t)<G(t-1)$, and $\cost(G)$ is the number of violated hops of
both heads (counted per head) plus the number $\desc(G)$ of descents. The \emph{error} of a hop is
$e(p,u)=(G(u)-G(p)-1)\bmod q\in\{0,\dots,q-1\}$.

\begin{lemma}[Lemma Zq of the fourth paper]\label{lem:Zq4}
\cited{the fourth paper~\cite[Lemma~5.15]{FourthPaper2026}; Lean-verified there (\lean{MotionDepthCodes.lemma_Zq},
\lean{MotionDepthCodes.card_descents_eq}).}
On the $r$-fold cover, every label satisfies $q\cdot\cost(G)=n_{\mathrm h}+\sum_{\text{hops}}e+q\cdot\desc(G)$ and
$q\cdot\desc(G)=\sum_{t\in\Z_T}\chi_t$, where $\chi_t=(G(t)-G(t-1))\bmod q$.
\end{lemma}

The \emph{level} of a hop is the least level of a turn of its own head strictly inside it (inside any lift to~$\Z$); the
owner is such a turn, so levels are defined. The hops of head~$h$ of level~$i$ form the \emph{class} $(h,i)$. A label
\emph{$\theta$-serves} level~$i$ for head~$h$ if the class $(h,i)$ has total error at most
$\theta qn_{\mathrm h}/(2(m+1))$, and \emph{doubly $\theta$-serves} it if it $\theta$-serves it for both heads
\cite[Section~5.7]{FourthPaper2026}. By \cref{lem:dnp-classes} below, every class has exactly $n_{\mathrm h}/(2(m+1))$
hops, so this says that the mean error of the class is at most~$\theta q$.

\begin{conjecture}[Conjecture A$'$ of the fourth paper]\label{conj:Aprime}
There are $\theta_0>0$ and constants $C,C'$, and for all $s_0$ and~$m$ an odd shift~$\xi$, such that for all $r$ and~$q$
every label~$G$ with $\cost(G)\le n_{\mathrm h}/8$ on the $r$-fold cover of the pair $x_1=U(s_0,m,1)$,
$x_2=x_1(\cdot+\xi)$, has at most $C\log_2q+C'$ doubly $\theta_0$-served levels~\cite[Conjecture~5.25]{FourthPaper2026}.
\end{conjecture}

By the per-class split of Lemma~Zq~\cite[Corollary~5.16]{FourthPaper2026}, $\cost(G)$ is at least the total error of
all hops divided by~$q$, so a level that is not doubly $\theta_0$-served contributes more than
$\theta_0n_{\mathrm h}/(2(m+1))$ to $\cost(G)$; hence a label of cost at most $\eps n_{\mathrm h}$ with $\eps\le1/8$ and
$\eps<\theta_0/2$ doubly serves more than $(m+1)(1-2\eps/\theta_0)$ levels. Under
Conjecture~A$'$ this forces $q\ge\exp(c(1-2\eps/\theta_0)(m+1)-c')$ for constants $c,c'>0$, and the same statement on the
walk DAG of many periods would make the pairs depth-robust as $m\to\infty$, so that $\TCP(2)$ fails
\cite[Section~5.7]{FourthPaper2026}. The fourth paper's clock labels cut every pair at depth $\lceil12/\eps\rceil^{m+1}-1$
with $\eps T$ edges~\cite[Proposition~5.13]{FourthPaper2026}.

\subsection{Formal status of the results}\label{sec:prelim-lean}

Every numbered result names its formal status in one of four ways. A line
\begin{center}
{\small\textsf{Lean:} \texttt{declaration, \dots}\hfill[\texttt{file}]}
\end{center}
gives the Lean declarations of a fully formalized result. The declarations live in the namespace
\texttt{BoundedNesting} and the files in the directory \leanfile{lean/BoundedNesting/}. The text after such a statement,
labelled \emph{Outline of the formal proof}, is a guide to the formal proof; it is a convenience, not a proof. A line
beginning \textsf{Lean (partial)} marks a result whose formal counterpart covers all but an explicitly named part, which
the text proves; on a definition, it names the notions that have no Lean counterpart. The words \textsf{Not formalized} mark a result with no formal counterpart; it comes with a complete
proof, labelled \emph{Proof (not formalized)}. A line beginning \textsf{Cited} marks a result of an earlier paper of the
series that is used without proof, and says where it is proved and whether it is Lean-verified there. Where the English
statement simplifies, the Lean statement is authoritative. Results from the literature are cited. Numerical observations
appear only in remarks labelled as numerical, heuristics are labelled as such, and claims that have only a proof sketch are
stated in remarks beginning \emph{Proof sketch only}; none of them is used in proofs. Numbered definitions carry a status line
only when they have a Lean counterpart.

The library depends on Lean~4 and Mathlib (both v4.34.1) and on the Lean library of the fourth paper, which imports those
of the first two papers and the repository's shared library \texttt{lean-machine}, all through relative paths. It
contains no \texttt{sorry}, no \texttt{native\_decide} and no unsafe evaluation, and every declaration cited here depends
only on the axioms \texttt{propext}, \texttt{Classical.choice} and \texttt{Quot.sound} (\cref{app:verify}).
\Cref{app:index} lists every declaration of this library cited in the paper.
